%%%%%%%%%%%%%%%%
%%% gen 48 %%%
%%%%%%%%%%%%%%%%

\Chapter{48}

\Verse{1}
{
	\hE{וַיְהִי אַחֲרֵי הַדְּבָרִים הָאֵלֶּה וַיֹּאמֶר לְיוֹסֵף הִנֵּה אָבִיךָ חֹלֶה וַיִּקַּח אֶת שְׁנֵי בָנָיו עִמּוֹ אֶת מְנַשֶּׁה וְאֶת אֶפְרָיִם׃}
}
{
	\eE{
		And it was after these things, and one said to Joseph,
		“Here, your father is sick.” And he took his two sons with
		him, Man-asseh and Ephraim.
	}
}
{
}

\Verse{2}
{
	\hE{וַיַּגֵּד לְיַעֲקֹב וַיֹּאמֶר הִנֵּה בִּנְךָ יוֹסֵף בָּא אֵלֶיךָ וַיִּתְחַזֵּק יִשְׂרָאֵל וַיֵּשֶׁב עַל הַמִּטָּה׃}
}
{
	\eE{
		And one told Jacob and said, “Here, your son Joseph is
		coming to you.” And Israel fortified himself and sat up on
		the bed.
	}
}
{
}

\Verse{3}
{
	\hP{וַיֹּאמֶר יַעֲקֹב אֶל יוֹסֵף אֵל שַׁדַּי נִרְאָה אֵלַי בְּלוּז בְּאֶרֶץ כְּנָעַן וַיְבָרֶךְ אֹתִי׃}
}
{
	\eP{
		And Jacob said to Joseph, “El Shadday appeared to me in
		Luz in the land of Canaan, and He blessed me
	}
}
{
}

\Verse{4}
{
	\hP{וַיֹּאמֶר אֵלַי הִנְנִי מַפְרְךָ וְהִרְבִּיתִךָ וּנְתַתִּיךָ לִקְהַל עַמִּים וְנָתַתִּי אֶת הָאָרֶץ הַזֹּאת לְזַרְעֲךָ אַחֲרֶיךָ אֲחֻזַּת עוֹלָם׃}
}
{
	\eP{
		and said to me, ‘Here, I’m making you fruitful and
		multiplying you, and I’ll make you into a community of
		peoples, and I’ll give this land to your seed after you,
		an eternal possession.’
	}
}
{
%	\footnote{
%		I’m making you fruitful and multiplying you, and I’ll make you … But God
%		did not say, “I’ll make you fruitful…” God said, “Be fruitful …”
%		(35:11). Why does Jacob tell Joseph that God promised to do it when God
%		actually told him to do it? Perhaps it is because Jacob has only one
%		more son (Benjamin) after God tells him this, and so the becoming
%		fruitful must refer to the births of his grandchildren and
%		great-grandchildren. Jacob would not see this as in his power but rather
%		as God’s doing, and so he understands God’s words not as a command but
%		as a promise and a blessing.
%	}
}

\Verse{5}
{
	\hP{וְעַתָּה שְׁנֵי בָנֶיךָ הַנּוֹלָדִים לְךָ בְּאֶרֶץ מִצְרַיִם עַד בֹּאִי אֵלֶיךָ מִצְרַיְמָה לִי הֵם אֶפְרַיִם וּמְנַשֶּׁה כִּרְאוּבֵן וְשִׁמְעוֹן יִהְיוּ לִי׃}
}
{
	\eP{
		And now, your two sons who were born to you in the land of
		Egypt by my arrival to you at Egypt: they’re mine. Ephraim
		and Man-asseh will be like Reuben and Simeon to me.
	}
}
{
}

\Verse{6}
{
	\hP{וּמוֹלַדְתְּךָ אֲשֶׁר הוֹלַדְתָּ אַחֲרֵיהֶם לְךָ יִהְיוּ עַל שֵׁם אֲחֵיהֶם יִקָּרְאוּ בְּנַחֲלָתָם׃}
}
{
	\eP{
		And your offspring that you’ll have after them will be
		yours. They shall be called by their brothers’ names with
		regard to their inheritance.
	}
}
{
}

\Verse{7}
{
	\hP{וַאֲנִי בְּבֹאִי מִפַּדָּן מֵתָה עָלַי רָחֵל בְּאֶרֶץ כְּנַעַן בַּדֶּרֶךְ בְּעוֹד כִּבְרַת אֶרֶץ לָבֹא אֶפְרָתָה וָאֶקְבְּרֶהָ שָּׁם בְּדֶרֶךְ אֶפְרָת הִוא בֵּית לָחֶם׃}
}
{
	\eP{
		And I: when I was coming from Paddan Aram, Rachel died by
		me in the land of Canaan on the way when there was still a
		span of land to come to Ephrat, and I buried her there on
		the Ephrat road. That’s Bethlehem.”
	}
}
{
}

\Verse{8}
{
	\hE{וַיַּרְא יִשְׂרָאֵל אֶת בְּנֵי יוֹסֵף וַיֹּאמֶר מִי אֵלֶּה׃}
}
{
	\eE{And Israel saw Joseph’s sons and said, “Who are these?”}
}
{
}

\Verse{9}
{
	\hE{וַיֹּאמֶר יוֹסֵף אֶל אָבִיו בָּנַי הֵם אֲשֶׁר נָתַן לִי אֱלֹהִים בָּזֶה וַיֹּאמַר קָחֶם נָא אֵלַי וַאֲבָרְכֵם׃}
}
{
	\eE{
		And Joseph said to his father, “They’re my sons, whom {\God}
		has given me here.” And he said, “Bring them to me, and
		I’ll bless them.”
	}
}
{
}

\Verse{10}
{
	\hE{וְעֵינֵי יִשְׂרָאֵל כָּבְדוּ מִזֹּקֶן לֹא יוּכַל לִרְאוֹת וַיַּגֵּשׁ אֹתָם אֵלָיו וַיִּשַּׁק לָהֶם וַיְחַבֵּק לָהֶם׃}
}
{
	\eE{
		And Israel’s eyes were heavy from old age. He was not able
		to see. And he brought them close to him, and he kissed
		them and embraced them.
	}
}
{
%	\footnote{
%		Israel’s eyes were heavy from old age. The analogy to his father is
%		obvious. When Isaac had been old, on his deathbed, unable to see, Jacob
%		had come and appropriated his brother Esau’s blessing. Now Jacob himself
%		is old, on his deathbed, unable to see, and he favors the younger son,
%		Ephraim, in his blessing.
%	}
}

\Verse{11}
{
	\hE{וַיֹּאמֶר יִשְׂרָאֵל אֶל יוֹסֵף רְאֹה פָנֶיךָ לֹא פִלָּלְתִּי וְהִנֵּה הֶרְאָה אֹתִי אֱלֹהִים גַּם אֶת זַרְעֶךָ׃}
}
{
	\eE{
		And Israel said to Joseph, “I didn’t expect to see your
		face; and, here, {\God} has shown me your seed as well.”
	}
}
{
}

\Verse{12}
{
	\hE{וַיּוֹצֵא יוֹסֵף אֹתָם מֵעִם בִּרְכָּיו וַיִּשְׁתַּחוּ לְאַפָּיו אָרְצָה׃}
}
{
	\eE{
		And Joseph brought them out from between his knees, and he
		bowed, his nose to the ground.
	}
}
{
}

\Verse{13}
{
	\hE{וַיִּקַּח יוֹסֵף אֶת שְׁנֵיהֶם אֶת אֶפְרַיִם בִּימִינוֹ מִשְּׂמֹאל יִשְׂרָאֵל וְאֶת מְנַשֶּׁה בִשְׂמֹאלוֹ מִימִין יִשְׂרָאֵל וַיַּגֵּשׁ אֵלָיו׃}
}
{
	\eE{
		And Joseph took the two of them, Ephraim in his right
		hand, at Israel’s left; and Manasseh in his left hand, at
		Israel’s right; and he brought them over to him.
	}
}
{
}

\Verse{14}
{
	\hE{וַיִּשְׁלַח יִשְׂרָאֵל אֶת יְמִינוֹ וַיָּשֶׁת עַל רֹאשׁ אֶפְרַיִם וְהוּא הַצָּעִיר וְאֶת שְׂמֹאלוֹ עַל רֹאשׁ מְנַשֶּׁה שִׂכֵּל אֶת יָדָיו כִּי מְנַשֶּׁה הַבְּכוֹר׃}
}
{
	\eE{
		And Israel put out his right hand and placed it on
		Ephraim’s head—and he was the younger—and his left hand on
		Manasseh’s head. He crossed his hands—because Manasseh was
		the firstborn.
	}
}
{
}

\Verse{15}
{
	\hE{וַיְבָרֶךְ אֶת יוֹסֵף וַיֹּאמַר הָאֱלֹהִים אֲשֶׁר הִתְהַלְּכוּ אֲבֹתַי לְפָנָיו אַבְרָהָם וְיִצְחָק הָאֱלֹהִים הָרֹעֶה אֹתִי מֵעוֹדִי עַד הַיּוֹם הַזֶּה׃}
}
{
	\eE{
		And he blessed Joseph and said, “The {\God} before whom my
		fathers, Abraham and Isaac, walked, who shepherded me from
		my start to this day,
	}
}
{
%	\footnote{
%		before whom my fathers walked. Note here and in Gen 6:9 and 17:1 that
%		this is technical covenant terminology known from ancient Near Eastern
%		documents to mean loyalty to one’s partner in a covenant.
%	}
}

\Verse{16}
{
	\hE{הַמַּלְאָךְ הַגֹּאֵל אֹתִי מִכּל רָע יְבָרֵךְ אֶת הַנְּעָרִים וְיִקָּרֵא בָהֶם שְׁמִי וְשֵׁם אֲבֹתַי אַבְרָהָם וְיִצְחָק וְיִדְגּוּ לָרֹב בְּקֶרֶב הָאָרֶץ׃}
}
{
	\eE{
		the angel who redeemed me from all bad, may He bless the
		boys. And may my name be called on them, and the name of
		my fathers, Abraham and Isaac. And may they spawn into a
		great number within the earth.”
	}
}
{
}

\Verse{17}
{
	\hE{וַיַּרְא יוֹסֵף כִּי יָשִׁית אָבִיו יַד יְמִינוֹ עַל רֹאשׁ אֶפְרַיִם וַיֵּרַע בְּעֵינָיו וַיִּתְמֹךְ יַד אָבִיו לְהָסִיר אֹתָהּ מֵעַל רֹאשׁ אֶפְרַיִם עַל רֹאשׁ מְנַשֶּׁה׃}
}
{
	\eE{
		And Joseph saw that his father had placed his right hand
		on Ephraim’s head, and it was bad in his eyes, and he held
		up his father’s hand to turn it from on Ephraim’s head
		onto Manasseh’s head,
	}
}
{
}

\Verse{18}
{
	\hE{וַיֹּאמֶר יוֹסֵף אֶל אָבִיו לֹא כֵן אָבִי כִּי זֶה הַבְּכֹר שִׂים יְמִינְךָ עַל רֹאשׁוֹ׃}
}
{
	\eE{
		and Joseph said to his father, “Not like that, my father,
		because this is the firstborn. Set your right hand on his
		head.”
	}
}
{
}

\Verse{19}
{
	\hE{וַיְמָאֵן אָבִיו וַיֹּאמֶר יָדַעְתִּי בְנִי יָדַעְתִּי גַּם הוּא יִהְיֶה לְּעָם וְגַם הוּא יִגְדָּל וְאוּלָם אָחִיו הַקָּטֹן יִגְדַּל מִמֶּנּוּ וְזַרְעוֹ יִהְיֶה מְלֹא הַגּוֹיִם׃}
}
{
	\eE{
		And his father refused and said, “I know, my son, I know.
		He, too, will become a people; and he, too, will be great.
		But in fact his little brother will be greater than he,
		and his seed will be full-fledged of nations.”
	}
}
{
%	\footnote{
%		full-fledged of nations. The unusual phrase ml’-haggyim is usually taken
%		to mean a multitude of nations, but that makes no particular sense in
%		terms of the fate of the single tribe of Ephraim. Elsewhere ml’ can mean
%		a full unit among a group (as in 2 Sam 8:2). It may mean here that
%		Ephraim will be thought of as a nation, for Ephraim later comes to
%		dominate the kingdom of Israel, and the name Ephraim is sometimes used
%		to refer to the entire Israelite kingdom (Isa 7:2–17; Hos 5:3; 6:10;
%		7:1).
%	}
}

\Verse{20}
{
	\hE{וַיְבָרְכֵם בַּיּוֹם הַהוּא לֵאמוֹר בְּךָ יְבָרֵךְ יִשְׂרָאֵל לֵאמֹר יְשִׂמְךָ אֱלֹהִים כְּאֶפְרַיִם וְכִמְנַשֶּׁה וַיָּשֶׂם אֶת אֶפְרַיִם לִפְנֵי מְנַשֶּׁה׃}
}
{
	\eE{
		And he blessed them in that day, saying, “Israel will
		bless with you, saying: ‘May {\God} make you like Ephraim and
		like Manasseh.’” And he set Ephraim before Manasseh.
	}
}
{
}

\Verse{21}
{
	\hE{וַיֹּאמֶר יִשְׂרָאֵל אֶל יוֹסֵף הִנֵּה אָנֹכִי מֵת וְהָיָה אֱלֹהִים עִמָּכֶם וְהֵשִׁיב אֶתְכֶם אֶל אֶרֶץ אֲבֹתֵיכֶם׃}
}
{
	\eE{
		And Israel said to Joseph, “Here, I’m dying. And {\God} will
		be with you and will bring you back to your fathers’ land.
		And I’ve given you one shoulder over your brothers, which
		I took from the Amorite’s hand with my sword and my bow.”
	}
}
{
%	\footnote{
%		one shoulder over your brothers. This expression appears to refer to
%		Joseph’s getting two tribes while each of his brothers gets only one.
%		But it also puns on the word for shoulder, Hebrew š kem (Shechem).
%		Shechem is the name of the city that will one day be the capital of the
%		kingdom of Israel, and it is located in one of the Joseph tribes
%		(Manasseh).
%	}
}

\Verse{22}
{
	\hE{וַאֲנִי נָתַתִּי לְךָ שְׁכֶם אַחַד עַל אַחֶיךָ אֲשֶׁר לָקַחְתִּי מִיַּד הָאֱמֹרִי בְּחַרְבִּי וּבְקַשְׁתִּי׃}
}
{
	\eE{}
}
{
}
